% Conteudo migrado de thesis/main.tex.
% O titulo do capitulo fica definido em thesis.tex.


Para que a solução sem fio seja viável em ambiente de competição, ela deve
satisfazer um conjunto de requisitos funcionais e não funcionais que
refletem tanto as restrições operacionais da equipe quanto as exigências
técnicas de integração com o sistema existente.

\section{Desempenho e Latência}

O sistema deve ser capaz de transmitir dados com latência compatível com as frequências
de aquisição já praticadas pela equipe. Atrasos na entrega dos dados
comprometem a sincronização temporal com os demais canais do sistema e podem
tornar a informação inutilizável para análise.

\section{Conexão Instantânea}

O sistema não pode depender de um processo de estabelecimento de conexão entre
os nós antes do início da transmissão. Em cenários como um sensor de
passagem externo ao veículo, o carro pode estar em alta velocidade quando passa
pelo ponto de medição — qualquer overhead de handshake inviabilizaria a coleta
do dado. A comunicação deve estar sempre pronta para transmitir, sem negociação
prévia entre os dispositivos.

\section{Múltiplas Conexões Simultâneas}

O sistema deve suportar múltiplos nós sensores transmitindo simultaneamente
para um ou mais receptores. Isso é especialmente relevante para a instrumentação
das quatro rodas do veículo, onde todos os sensores precisam operar em
paralelo e ter seus dados recebidos e injetados na rede CAN de forma
coordenada.

\section{Transparência com o Sistema CAN Existente}

A integração com o barramento CAN deve ser feita com o mínimo de modificações
possível na arquitetura existente. O sistema sem fio deve preservar os
identificadores CAN originais de cada mensagem e manter as verificações de
integridade nativas do protocolo — incluindo CRC e acknowledgment —
garantindo que o restante da rede trate os dados recebidos sem fio de forma
idêntica aos dados provenientes de nós cabeados.

\section{Requisitos Físicos e de Custo}

Os nós sensores devem ser fisicamente compactos e leves, uma vez que serão
instalados em regiões onde massa e volume são restrições relevantes — como
rodas e componentes aerodinâmicos. O custo total do sistema deve ser compatível
com o orçamento de uma equipe universitária, priorizando componentes
acessíveis sem comprometer os requisitos funcionais.

\section{Modularidade}

A arquitetura do sistema deve ser modular, permitindo adicionar novos nós
sensores sem alterações na infraestrutura existente. Cada novo ponto de medição
deve poder ser incorporado ao sistema de forma independente, com impacto
mínimo nos demais nós e sem necessidade de reconfiguração do receptor ou do
barramento CAN.
% what are the requirements for a successful system

